\section{Сборка и использование}
\label{sec:running}

\subsection{Зависимости}
Для приложения необходимы компилятор с поддержкой C++17, \code{make},
\code{qmake} и библиотеки Qt Core, Gui, Widgets. Файл проекта
\code{PharmDelivery.pro} задаёт \code{CONFIG += c++17 cmdline}
и подключает каталоги корня проекта и \code{back/} к поиску заголовков.
Исходный набор модулей Qt указан как \code{QT = core widgets}; Gui подключается
через зависимости Widgets.

Для модульных тестов достаточно C++17-компилятора и Bash; они не используют
\code{main.cpp} и собираются без Qt. Проверки командной строки требуют
собранного приложения и Python 3 со стандартной библиотекой.

\subsection{Сборка приложения}
Выполните из корня проекта:
\begin{lstlisting}[language=bash]
cd /home/toster/prac-backend
mkdir -p build-backend
cd build-backend
qmake6 ../PharmDelivery.pro
make -j2
./PharmDelivery --help
\end{lstlisting}
Для Qt 5 замените \code{qmake6} на \code{qmake}. Ожидаемое имя исполняемого
файла --- \code{PharmDelivery}. Отдельный каталог сборки сохраняет корень
исходников свободным от объектных файлов.

После разбора аргументов \code{main.cpp} задаёт \code{QT_QPA_PLATFORM=offscreen},
если эта переменная не задана или пуста. Это позволяет выполнять модель без
графической сессии при наличии платформенного плагина Qt. Уже заданное значение
переменной сохраняется.

\subsection{Параметры модели}
Аргумент и его значение передаются отдельными элементами командной строки,
например \code{--days 15}. Числа с дробной частью в примерах используют точку.
Все вещественные параметры должны быть конечными.

{\small
\begin{longtable}{@{}L{0.30\textwidth}L{0.12\textwidth}L{0.49\textwidth}@{}}
\toprule
Параметр & По умолчанию & Значение и ограничения \\
\midrule
\endhead
\code{--days}, \code{-N} & 15 & Число дней, целое от 10 до 25. \\
\code{--couriers}, \code{-M} & 5 & Число курьеров, целое от 3 до 9. \\
\code{--medicines}, \code{-K} & 25 & Размер каталога, целое от 15 до 35. \\
\code{--markup} & 25 & Наценка в процентах, от 0 до 10000. \\
\code{--card-discount} & 5 & Скидка по карте в процентах, от 0 до 100. \\
\code{--bulk-discount} & 3 & Скидка без карты при сумме строго больше 1000 рублей,
в процентах, от 0 до 100. \\
\code{--regular-discount} & 5 & Дополнительная скидка постоянному покупателю,
в процентах, от 0 до 100. \\
\code{--max-discount} & 9 & Верхний предел итоговой скидки в процентах, от 0 до 9. \\
\code{--initial-count} & 80 & Начальные упаковки каждого лекарства,
целое от 0 до 1000000; при \code{--stock} партии берутся из файла. \\
\code{--regular-customers} & 10 & Число постоянных покупателей,
целое от 0 до 100000. \\
\code{--seed} & 5489 & Начальное состояние генератора,
целое от 0 до 4294967295. \\
\code{--stock FILE} & Не задан & Файл начальных партий в формате CSV.
Заменяет весь автоматически создаваемый запас. \\
\bottomrule
\end{longtable}
}

Размеры отдельных скидок могут превышать общий предел: при расчёте цены
сумма всё равно ограничивается \code{--max-discount}. Валидация
\code{SimulationParameters} выполняется до инициализации модели.

\subsection{Параметры вывода}
\begin{tabularx}{\textwidth}{@{}L{0.30\textwidth}Y@{}}
\toprule
Параметр & Поведение \\
\midrule
\code{--help}, \code{-h} & Показать справку и завершить приложение. \\
\code{--quiet} & Показать только итоговую статистику. \\
\code{--orders} & Добавить сведения о каждой доставке к дневным отчётам. \\
\code{--report FILE} & Записать полный JSON-отчёт после последнего дня. \\
\bottomrule
\end{tabularx}

Без флагов вывода приложение печатает каталог, события каждого дня, остатки,
нагрузку курьеров и итоги. Совместное применение \code{--quiet} и \code{--orders}
оставляет только итоги: дневной вывод при \code{--quiet} отключён.
Родительская директория файла \code{--report} должна существовать.
Существующий файл отчёта открывается с перезаписью.

\subsection{Примеры экспериментов}
Базовый запуск из каталога сборки:
\begin{lstlisting}[language=bash]
./PharmDelivery --days 15 --couriers 5 --medicines 25 \
  --markup 25 --seed 42
\end{lstlisting}

Эксперимент со скидками и сохранением результата:
\begin{lstlisting}[language=bash]
./PharmDelivery --days 25 --couriers 7 --medicines 35 \
  --markup 20 --card-discount 4 --bulk-discount 2 \
  --regular-discount 3 --max-discount 9 \
  --seed 71 --quiet --report experiment.json
\end{lstlisting}

Использование начального склада из проекта:
\begin{lstlisting}[language=bash]
./PharmDelivery --days 10 --couriers 5 --medicines 15 \
  --stock ../examples/initial-stock.csv \
  --seed 42 --report stock-experiment.json
\end{lstlisting}

Один и тот же \code{seed} и полный набор параметров воспроизводят эксперимент
при неизменном файле начального склада, коде и реализации стандартной библиотеки.
Изменение параметров может менять число обращений к генератору: одинаковый
\code{seed} не означает идентичный поток заказов во всех сравниваемых сценариях.

\subsection{Ошибки и завершение}
Успешный запуск и вывод справки завершаются с кодом 0. Исключение, перехваченное
в \code{main.cpp}, приводит к сообщению с префиксом «Ошибка:» в стандартном
потоке ошибок и коду 1. Типичные причины: неизвестный аргумент, пропущенное
значение, выход параметра за диапазон, некорректный CSV или недоступный путь
отчёта. Ошибки среды Qt, возникающие при создании \code{QApplication}, могут
завершать процесс через механизм Qt и требуют проверки установки плагинов.
